\section{Shrinking viability and coherent approximate names}
\label{sec:shrinking}

\subsection{The remaining coherence problem}

Finite amalgamation solves every theorem with a fixed finite horizon.  It does
not by itself produce one infinite coherent branch.  Nor should it: an
infinite finitely branching tree can have paths of every finite length while a
uniform choice of an infinite branch may require additional logical or
computational strength, depending on how the tree is presented.

Most mathematical applications need less.  They do not need a stable code for
the same branch at every level; they need observations that become mutually
close.  If every full finite solution has nearly the same observation at the
requested level, then arbitrary finite choices can be determinized without
choosing an infinite branch.  This is the content of the shrinking-viability
theorem.

\subsection{Metric observations}

For each good candidate \(a<B(p,k)\), let
\(\Obs(p,k,a)\) be a rational or coded metric observation.  Let
\(d(x,y)\) be a primitive-recursive rational-valued pseudometric on such
observations.  We require the ordinary finite metric laws as embedded
inequalities.  Exact equality of codes is not required: different candidates
may denote the same approximate object.

Assume a primitive-recursive edge budget \(\delta(p,k)\ge0\) such that
\begin{equation}
 \Good(p,k,a)\wedge\Good(p,k+1,b)\wedge\Compat(p,k,a,b)
 \quad\Longrightarrow\quad
 d(\Obs(p,k,a),\Obs(p,k+1,b))\le\delta(p,k).
 \label{eq:edge-budget}
\end{equation}
Let
\[
 \Tail(p,k,\ell)=\sum_{j=k}^{\ell-1}\delta(p,j)
 \qquad(k\le\ell)
\]
be the finite telescoping budget.

Choose a nondecreasing primitive-recursive horizon function \(H(p,k)\) with
\(H(p,k)\ge k\).  Define the finite viable cloud
\[
 V_{p,k}:=\{a<B(p,k):\FullV(p,H(p,k),k,a)=1\}.
\]
It is a finite, decidable set.  Its canonical representative is
\[
 s(p,k)=\min V_{p,k},
\]
with any fixed default value when the set is empty.  In applications, a
separate hypothesis proves nonemptiness.

\subsection{The shrinking-viability theorem}

\begin{theorem}[Shrinking viability]
\label{thm:shrinking-viability}
Suppose iPRA proves the following numerical contracts for every admissible
parameter \(p\).

\begin{enumerate}[label=(SV\arabic*)]
\item \emph{Nonempty finite solutions:}
      \(V_{p,k}\ne\varnothing\) for every \(k\).
\item \emph{Monotone horizons:}
      if \(k\le\ell\), then \(H(p,k)\le H(p,\ell)\).
\item \emph{Horizontal shrinking:} for all \(a,b\in V_{p,k}\),
      \[
       d(\Obs(p,k,a),\Obs(p,k,b))\le\varepsilon(p,k).
      \]
\item \emph{Edge control:} every compatible edge satisfies
      \eqref{eq:edge-budget}.
\end{enumerate}
Then the primitive-recursive selector
\[
 x(p,k):=\Obs(p,k,s(p,k))
\]
satisfies, for all \(k\le\ell\),
\begin{equation}
 d(x(p,k),x(p,\ell))
 \le \varepsilon(p,k)+\Tail(p,k,\ell).
 \label{eq:shrinking-bound}
\end{equation}
Consequently, if a primitive-recursive modulus \(r(p,n)\) satisfies
\[
 k\ge r(p,n)\to
 \varepsilon(p,k)+\sum_{j=k}^{m-1}\delta(p,j)\le 2^{-n}
 \quad\text{for every }m\ge k,
\]
then \(k\mapsto x(p,r(p,k))\) is a regular Cauchy name with an explicit
primitive-recursive rate.
\end{theorem}

\begin{proof}
Fix \(p\) and \(k\le\ell\).  By (SV1), the selected candidate
\(s(p,\ell)\) lies on a complete good compatible path
\[
 c_0,c_1,\ldots,c_{H(p,\ell)}
 \qquad c_\ell=s(p,\ell).
\]
Here and below the path is obtained by bounded search from the Boolean
\(\FullV\) assertion; no unbounded choice is used.

Let \(b=c_k\).  Since horizons are nondecreasing,
\(H(p,k)\le H(p,\ell)\).  Truncating the displayed path at level \(H(p,k)\)
shows that \(b\in V_{p,k}\).  The selected candidate \(s(p,k)\) also lies in
\(V_{p,k}\), so horizontal shrinking gives
\[
 d(\Obs(p,k,s(p,k)),\Obs(p,k,b))\le\varepsilon(p,k).
\]
Along the segment \(c_k,c_{k+1},\ldots,c_\ell\), edge control and repeated
triangle inequalities give
\[
 d(\Obs(p,k,b),\Obs(p,\ell,s(p,\ell)))
 \le \sum_{j=k}^{\ell-1}\delta(p,j).
\]
One final triangle inequality proves \eqref{eq:shrinking-bound}.

Every operation in this argument is finite.  In the object theory, the path
is the least bounded path certified by \(\FullV\); membership in the two
viable clouds is a Boolean equation; the telescoping inequality is proved by
quantifier-free induction on \(\ell-k\).  Universal introduction over
\(k,\ell\) is performed after that induction.  The final regularity claim is
obtained by substituting the supplied modulus.
\end{proof}

\subsection{Where horizontal shrinking comes from}

The hypothesis that a viable cloud has small diameter is quantitative
uniqueness, not an unexplained compactness principle.  Several standard
sources are reusable.

\paragraph{Isolating intervals or balls.}
If every full finite solution at level \(k\) is certified to lie in the same
rational ball of radius \(\rho_k\), then
\(\varepsilon_k=2\rho_k\).

\paragraph{Strong convexity.}
Suppose candidates satisfy
\(f(a)\le m_k+\eta_k\), while
\(f(x)\ge f(x_*)+\lambda d(x,x_*)^2\).  Then every viable candidate is within
\(\sqrt{\eta_k/\lambda}\) of the unique minimizer, so the cloud diameter is at
most twice that bound.  In a strictly finitary development, the square-root
bound is replaced by a rational dyadic upper bound verified by squaring.

\paragraph{Residual-to-distance estimates.}
If a certified inequality
\[
 d(x,Z)\le C\,R(x)
\]
relates residual to distance from one isolated solution or one small cluster,
and every viable candidate has residual at most \(\eta_k\), then the cloud
diameter is bounded by twice \(C\eta_k\).

\paragraph{Algebraic separation.}
For a polynomial with exact algebraic coefficient data, isolating disks and a
root-separation bound group approximations into disjoint clusters.  Within a
cluster the viable cloud shrinks; at collisions the cluster or root multiset,
not an ordered root, is the stable object.

\paragraph{Contraction certificates.}
If a rational box is mapped into itself by a contraction of factor
\(q<1\), any two candidates surviving \(n\) refinement steps differ by at
most \(q^n\) times the initial diameter plus the verified rounding budget.

These mechanisms are exactly the data absent from a bare qualitative
uniqueness statement.  The theorem refuses to infer a modulus from uniqueness
alone.

\subsection{Nested balls as a special case}

\begin{corollary}[Nested-ball name]
\label{cor:nested-ball-name}
Suppose a certificate at level \(k\) is a rational center \(c_k\) and radius
\(r_k\), every full path chooses nested balls, and every viable ball at level
\(k\) intersects one common rational ball of radius \(r_k\).  If
\(r_k\le 2^{-k-3}\) and compatible centers move by at most
\(2^{-k-3}\), then the least full-viable center, shifted by a fixed number of
precision levels, is a regular name.
\end{corollary}

\begin{proof}
Intersecting balls of radius \(r_k\) have centers at distance at most
\(2r_k\), so take \(\varepsilon_k=2r_k\).  The edge budgets form a geometric
series.  Apply \cref{thm:shrinking-viability} and absorb the constant factor by
shifting the requested precision.
\end{proof}

\subsection{Exact choice under a positive gap}

Shrinking gives approximate coherence.  Exact labels can be recovered only
when the labels are separated.

\begin{corollary}[Eventual exact label stability]
\label{cor:exact-label-gap}
Suppose candidates carry labels in a finite set, candidates with different
labels have observational distance at least \(g>0\), and the right side of
\eqref{eq:shrinking-bound} is eventually less than \(g\).  Then the selected
labels are eventually constant.  A primitive-recursive search for the first
precision where the certified error is below \(g\) yields an exact label.
\end{corollary}

\begin{proof}
Two selected candidates at sufficiently high precisions are closer than
\(g\); therefore they cannot have different labels.  The threshold is found
from the supplied numerical modulus, not from a qualitative assertion that a
gap exists.
\end{proof}

This is the correct interface for root ordering, pivot choice, chart choice,
sign, and projective orientation: approximate or relational data are
unconditional; exact labels require a positive rational separation witness.

\subsection{What happens without shrinking}

If (SV3) is unavailable, the framework still yields strong conclusions.
For every requested finite collection of precisions, finite amalgamation
produces a compatible family.  Downstream theorems that are stable under all
such families can consume the relation directly.  Examples include:

\begin{itemize}
\item existence of a root cluster meeting a residual bound;
\item the minimum value of a function without a selected minimizer;
\item approximate rotation data without a convergent axis;
\item a finite cover or parity witness in topology;
\item all representatives passing a finite test suite.
\end{itemize}

The absence of a coherent selector is then a property of the statement, not a
failure of the encoding.  This distinction prevents hidden uses of choice,
LLPO, or discontinuous branch selection.
